\subsection{Обратимость линейного, ограниченного снизу \texorpdfstring{\((\|Ax\|\geqslant k\|x\|)\)}{(||Ax|| >= k||x||)} оператора}

\begin{definition}[Обратимость оператора на образе]
	Пусть \(E_1,E_2\) --- нормированные пространства,
	\[
	A\in\mathcal{L}(E_1,E_2).
	\]
	Оператор \(A\) называется \textbf{обратимым на своём образе}, если для любого
	\[
	y\in \operatorname{Im}A
	\]
	уравнение
	\[
	Ax=y
	\]
	имеет единственное решение \(x\in E_1\).
	
	В этом случае определён обратный оператор
	\[
	A^{-1}:\operatorname{Im}A\to E_1,
	\]
	задаваемый правилом
	\[
	A^{-1}y=x
	\quad\Longleftrightarrow\quad
	Ax=y.
	\]
\end{definition}

\begin{remark}[Левый и правый обратные]
	Для оператора
	\[
	A:E_1\to E_2
	\]
	можно различать левый и правый обратные операторы.
	
	\begin{itemize}
		\item \textbf{Левый обратный} отвечает за единственность решения:
		\[
		LA=I_{E_1}.
		\]
		
		\item \textbf{Правый обратный} отвечает за существование решения на образе:
		\[
		AR=I_{\operatorname{Im}A}.
		\]
	\end{itemize}
	
	Если оба оператора существуют и совпадают, то получаем обычный обратный оператор
	\[
	A^{-1}.
	\]
	В этом билете обратный оператор рассматривается именно как оператор
	\[
	A^{-1}:\operatorname{Im}A\to E_1.
	\]
\end{remark}

\begin{definition}[Ограниченный снизу оператор]
	Пусть \(E_1,E_2\) --- нормированные пространства,
	\[
	A\in\mathcal{L}(E_1,E_2).
	\]
	Оператор \(A\) называется \textbf{ограниченным снизу}, если существует число
	\[
	k>0,
	\]
	такое что
	\[
	\|Ax\|_{E_2}\geqslant k\|x\|_{E_1}
	\qquad
	\forall x\in E_1.
	\]
\end{definition}

\begin{lemma}[Ограниченность снизу влечёт инъективность]
	Если оператор
	\[
	A\in\mathcal{L}(E_1,E_2)
	\]
	ограничен снизу, то
	\[
	\ker A=\{0\}.
	\]
	Следовательно, \(A\) инъективен.
\end{lemma}

\begin{proof}
	Пусть
	\[
	x\in\ker A.
	\]
	Тогда \(Ax=0\). Так как \(A\) ограничен снизу, существует \(k>0\), такое что
	\[
	\|Ax\|\geqslant k\|x\|.
	\]
	Но \(Ax=0\), поэтому
	\[
	0=\|Ax\|\geqslant k\|x\|.
	\]
	Так как \(k>0\), получаем
	\[
	\|x\|=0.
	\]
	Значит,
	\[
	x=0.
	\]
	Следовательно,
	\[
	\ker A=\{0\}.
	\]
\end{proof}

\begin{theorem}[Критерий ограниченности обратного оператора]
	Пусть \(E_1,E_2\) --- нормированные пространства,
	\[
	A\in\mathcal{L}(E_1,E_2).
	\]
	Тогда обратный оператор
	\[
	A^{-1}:\operatorname{Im}A\to E_1
	\]
	существует и ограничен тогда и только тогда, когда оператор \(A\) ограничен снизу, то есть
	\[
	\exists k>0
	\quad
	\forall x\in E_1
	\qquad
	\|Ax\|\geqslant k\|x\|.
	\]
\end{theorem}

\begin{proof}
	Докажем обе импликации.
	
	\textbf{1. Если \(A^{-1}\) существует и ограничен, то \(A\) ограничен снизу.}
	
	Пусть
	\[
	A^{-1}\in\mathcal{L}(\operatorname{Im}A,E_1).
	\]
	Возьмём произвольный элемент \(x\in E_1\). Тогда \(Ax\in\operatorname{Im}A\), и по определению
	обратного оператора
	\[
	A^{-1}Ax=x.
	\]
	Следовательно,
	\[
	\|x\|
	=
	\|A^{-1}Ax\|
	\leqslant
	\|A^{-1}\|\cdot\|Ax\|.
	\]
	Отсюда
	\[
	\|Ax\|
	\geqslant
	\frac{1}{\|A^{-1}\|}\|x\|.
	\]
	Значит, \(A\) ограничен снизу с константой
	\[
	k=\frac{1}{\|A^{-1}\|}.
	\]
	
	\textbf{2. Если \(A\) ограничен снизу, то \(A^{-1}\) существует и ограничен.}
	
	Пусть существует \(k>0\), такое что
	\[
	\|Ax\|\geqslant k\|x\|
	\qquad
	\forall x\in E_1.
	\]
	Докажем, что
	\[
	A^{-1}:\operatorname{Im}A\to E_1
	\]
	корректно определён и ограничен.
	
	\begin{itemize}
		\item \textbf{Существование обратного оператора на образе.}
		
		По лемме из ограниченности снизу следует
		\[
		\ker A=\{0\}.
		\]
		Значит, \(A\) инъективен.
		
		Поэтому для каждого \(y\in\operatorname{Im}A\) существует единственный \(x\in E_1\), такой что
		\[
		Ax=y.
		\]
		Следовательно, корректно определён оператор
		\[
		A^{-1}:\operatorname{Im}A\to E_1.
		\]
		
		\item \textbf{Линейность обратного оператора.}
		
		Пусть
		\[
		y_1,y_2\in\operatorname{Im}A,
		\qquad
		\alpha,\beta\in\K.
		\]
		Выберем \(x_1,x_2\in E_1\), такие что
		\[
		Ax_1=y_1,
		\qquad
		Ax_2=y_2.
		\]
		Тогда
		\[
		A(\alpha x_1+\beta x_2)
		=
		\alpha Ax_1+\beta Ax_2
		=
		\alpha y_1+\beta y_2.
		\]
		По единственности прообраза
		\[
		A^{-1}(\alpha y_1+\beta y_2)
		=
		\alpha x_1+\beta x_2
		=
		\alpha A^{-1}y_1+\beta A^{-1}y_2.
		\]
		Значит, \(A^{-1}\) линеен.
		
		\item \textbf{Ограниченность обратного оператора.}
		
		Возьмём произвольный
		\[
		y\in\operatorname{Im}A.
		\]
		Тогда \(y=Ax\) для некоторого \(x\in E_1\), причём
		\[
		A^{-1}y=x.
		\]
		Используем ограниченность \(A\) снизу:
		\[
		\|y\|
		=
		\|Ax\|
		\geqslant
		k\|x\|
		=
		k\|A^{-1}y\|.
		\]
		Отсюда
		\[
		\|A^{-1}y\|
		\leqslant
		\frac{1}{k}\|y\|.
		\]
	\end{itemize}
	
	Итак,
	\[
	A^{-1}\in\mathcal{L}(\operatorname{Im}A,E_1),
	\qquad
	\|A^{-1}\|\leqslant \frac{1}{k}.
	\]
\end{proof}

\begin{remark}[Полезное следствие: замкнутость образа]
	Если \(E_1\) --- банахово пространство и оператор
	\[
	A\in\mathcal{L}(E_1,E_2)
	\]
	ограничен снизу, то
	\[
	\operatorname{Im}A
	\]
	замкнут в \(E_2\).
	
	Действительно, пусть
	\[
	y_n\in\operatorname{Im}A,
	\qquad
	y_n\to y.
	\]
	Представим
	\[
	y_n=Az_n.
	\]
	Так как \(\{y_n\}\) сходится, она фундаментальна. По оценке снизу
	\[
	\|y_n-y_m\|
	=
	\|Az_n-Az_m\|
	=
	\|A(z_n-z_m)\|
	\geqslant
	k\|z_n-z_m\|.
	\]
	Значит, \(\{z_n\}\) фундаментальна. Так как \(E_1\) банахово, существует \(z\in E_1\), такое что
	\[
	z_n\to z.
	\]
	По непрерывности \(A\)
	\[
	Az_n\to Az.
	\]
	Но \(Az_n=y_n\to y\), поэтому
	\[
	y=Az\in\operatorname{Im}A.
	\]
	Следовательно, \(\operatorname{Im}A\) замкнут.
\end{remark}
